\subsubsection{端到端吞吐量与 Batch Size 扩展性}

为评估不同 KV Cache 压缩配置在长上下文解码场景下的端到端性能，本文在 Meta-Llama-3-8B-Instruct 上进行了 batch size 扫描实验。实验设置输入长度为 4096，输出长度为 256，batch size 取 1、2、3、4、5 和 6。对比配置包括 Dense 基线、50\% 稀疏、70\% 稀疏、50\% 稀疏 + 2bit 量化以及 70\% 稀疏 + 2bit 量化。量化版本采用本文优化后的计算核心，其中 Key 路径与 Value 路径分别调优，Value 路径在长上下文 decode 场景下启用了更优的 Split-K 与 tile 配置。

从端到端吞吐量结果可以看出，两条量化稀疏曲线在 batch size 为 1--6 的范围内均稳定超过对应的纯稀疏版本。以 70\% 稀疏度为例，量化后吞吐量由 22.35、36.31、53.01、60.21、67.20 和 67.37 tokens/s 提升到 23.10、40.30、53.67、60.41、70.60 和 74.39 tokens/s；以 50\% 稀疏度为例，量化后吞吐量由 22.11、35.16、48.26、56.37、61.45 和 61.80 tokens/s 提升到 22.44、38.32、55.05、62.61、67.91 和 72.03 tokens/s。该结果说明，在当前长上下文设置下，量化稀疏实现已经能够将访存压缩收益转化为稳定的端到端吞吐提升。

进一步比较 50\% 与 70\% 两组量化配置可知，70\% 稀疏 + 2bit 量化在 batch size 为 1、2、5 和 6 时取得最高吞吐量，而 50\% 稀疏 + 2bit 量化在 batch size 为 3 和 4 时略高于 70\% 稀疏 + 2bit 量化。这表明在中等 batch size 区间，较高稀疏度带来的访存收益与量化元数据、反量化开销之间仍存在权衡；而当 batch size 继续增大后，70\% 稀疏配置的压缩收益再次占优。

TTFT 与 TPOT 的结果同样反映了这一趋势。总体上，量化稀疏配置在中高 batch size 下能够显著降低首 token 之后的稳态解码开销。例如在 batch size 为 6 时，70\% 稀疏配置的 TPOT 为 62.59 ms/token，而 70\% 稀疏 + 2bit 量化的 TPOT 降低到 56.06 ms/token；50\% 稀疏配置的 TPOT 为 62.59 ms/token 的同量级，而 50\% 稀疏 + 2bit 量化为 59.08 ms/token。与此同时，量化版本在多数 batch size 上也保持了更低或接近的 TTFT，说明其优化并未以显著拉高首 token 延迟为代价。

需要说明的是，吞吐量统计基于完整 256-token 生成过程计算；考虑到单卡 24GB 显存限制，TTFT 与 TPOT 采用前 64 个 decode steps 的独立计时窗口进行估计。对于 Dense 基线以及部分纯稀疏配置，较大的 batch size 会在完整的 token-level 计时过程中触发显存限制，因此表中相应位置记为 ``--''。即便如此，从端到端吞吐量结果仍可清晰看出：在当前实验设置下，优化后的量化稀疏计算核心已经在 batch size 为 1--6 的范围内稳定优于对应的纯稀疏实现。

图~\ref{fig:end2end-throughput-bs} 与表~\ref{tab:end2end-throughput-bs}--表~\ref{tab:end2end-tpot-bs} 给出了完整结果。

\begin{figure}[t]
    \centering
    \includegraphics[width=\linewidth]{figures/throughput_vs_bs_five_configs_quant_tuned_v3.pdf}
    \caption{Llama-3 8B 在不同 batch size 下的端到端吞吐量对比（input=4096, output=256）。}
    \label{fig:end2end-throughput-bs}
\end{figure}
